\chapter{Boot2Root and Repair: VulnHub Basic Pentesting 1}

In this post, we'll cover a walkthrough of a VulnHub VM, "Basic Pentesting 1" by Josiah Pierce. \[[3]\]  Put together for his colleagues at Virginia Tech's computer security courses, this beginner's-level VM has some ready-to-exploit vulnerabilities that we worked on.  Later in this report, once inside the system, we'll look around, break some stuff, and fix the vulnerabilities that helped us get in.  Along the way, we'll look at some password cracking with \texttt{hashcat} and what it takes to stop time. 


\subsubsection{Spoilers in a walkthrough}
    Throughout this post, we'll reveal some aspects of the "Basic Pentesting 1" target.  Please note that by looking at some of these details, the VM's instructional quality as a target of unknown composition might be diminished to hackers who will want use the target for technical growth and professional development.  Experience recommends thorough use of the target VM before referring to resources such as these.

\subsection{Some tips on learning from VMs}
At a recent event, it was recommended to me that if you are hacking on a VM and get stuck, you can phone-a-friend.  One technique is to contact someone else who hacks, get them to do a minute of research, and then have them send you a clue.  That way, you can get a little help without having to totally cave in and read a walkthrough before you're ready.

When this was pointed out, the pentester noted that he'd learned a lot from working the walkthroughs.  One by one, he'd knock them out.  By reviewing directions and typing in the commands, he found himself being able to absorb the techniques and information.  It reminded me a lot of my days as a young programmer.  I'd go to the grocery store, buy a magazine, and then begin typing in the programs.  Since I could not afford to purchase commercial ones, if I wanted my computer to do anything, I'd have to type it in myself.  It's good to see that hackers today are learning this same way.  

\section{Planning}
In this section of the report, we'll cover setup, machine configuration, and basic communication checks.  We'll use NIST's four phases of pentesting \[[2]\] from 800-115 combined with modified Lockheed-Martin kill chain phase descriptions \[[1]\] to summarize key points of our findings in tables.

Published in Dec 2017, the VM is based on a Ubuntu distribution.  Screenshots show a "Marlinspike" admin account login.  We're not shown much else.  We discovered later that the VM will not reset itself on restart; so, keep a copy of the downloaded file handy.  

\subsection{Setup}

We used pretty much the same setup; except, since my attack system was getting out of date, I spun up a fresh copy of Kali on a VM to help with the \texttt{hashcat}.

\subsection{Summary properties of the machines}
One box is holding the VulnHub VM; it's running VirtualBox; we conducted some simple  \texttt{ping} checks to make sure that it could communicate with other machines on the SOHO network.  This target box is the VM running on a WIN10 laptop.  On the attacker box, we are using FreeBSD 10.3 and a VirtualBox hypervisor running Kali in a VM.

\begin{table}[H]
\caption{Initial Commo Checks and Basic Configuation}
\centering
\begin{tabular}{|l|l|l|}
\hline
Property & Value & Logical Impact \\
\hline
Target Machine:  VirtualBox VM & VulnHub Basic Pentesting 1 in VM & Target of unknown composition \\
Target Machine:  Host OS & WIN 10 & Target platform host OS \\
Attacker Machine:  VirtualBox VM & Kali & Attack platform inside VM \\
Attacker Machine:  Host OS & FreeBSD 10.3 RELEASE & Attack platform host OS \\
Target & 192.168.1.17 & \texttt{ping} OK \\
Attacker & 192.168.1.10 & \texttt{ping} OK \\
\hline
\end{tabular}
\end{table}

\subsection{Proving COMMO}
Since we don't have a need for stealth and do have direct access to the controls of both VMs, we started off with a simple commo check between the boxes.  By using:

\begin{lstlisting}
ping -c 5 <DESTINATION IP ADDRESS>
\end{lstlisting}


\section{Discovery and attack}

\subsection{Vulnerability scanning}
This time we skipped the Nessus scan and went with a simpler approach.  We opened with \texttt{nikto} scans to determine directory structures and \texttt{nmap} scans to find out about open ports and program versions.  With some vulnerabilities quickly found, we went ahead and looked those up.

\subsubsection{`nmap` and `nikto` scans}
These revealed to us that port 80 was open with Apache 2.4.18 on Ubuntu Linux.  There was a directory marked \texttt{/secret} that looked like a WordPress blog.  There was also a \texttt{vtcsec/secret/index.php}  A quick manual call to the page's source code showed that the WordPress blog looked like it had a comments section and a 2017 theme.

\begin{lstlisting}
nmap -sV 192.168.1.17
\end{lstlisting}

\begin{table}[H]
\caption{Scan results}
\centering
\begin{tabular}{|l|l|l|}
\hline
Port & Protocol & Program \\
\hline
21 & ftp & ProFTPD 1.3.3c \\
22 & ssh & OpenSSH 7.2p2 \\
80 & http & Apache httpd 2.4.18 \\
\hline
\end{tabular}
\end{table}

One of the things missing from the scans was an open port for MySQL.  Having run WordPress setups many times, I noticed that this implied that the site might be running without a database backend.  This implied a PHP engine; and it also suggests that we would not be able to break in with SQLi through \texttt{sqlmap}.  The presence of a server-side language like PHP was good news, since it increased the chances that we would have a ready-made exploit on hand.

Curious about the other two programs, we looked those up.  We quickly found web pages describing vulnerabilities. \[[4]\] The ProFTPD vuln comes with its own Metasploit exploit.  We went with that right away.  All we had to do was to set the target IP and hit \texttt{exploit}.  About 15 minutes into the exercise, and we were in.  As soon as the exploit took, we had root. 


All of this unfolded so quickly, we turned our attention to better understanding actions on the objective; and, we would also be able to turn our attention to later repairs.

\subsubsection{Analysis Check}
\begin{table}[H]
\caption{Discovery and Attack Cycle 1}
\centering
\begin{tabular}{|l|l|}
\hline
Action & Observation \\
\hline
nmap & WordPress and PHP likely available. \\
nikto & Open ports, program names and versions \\
Internet research of programs & Metasploit \texttt{exploit/unix/ftp/proftpd\_133c\_backdoor} \\
\hline
\end{tabular}
\end{table}

\begin{table}[H]
\caption{Command to Kill Chain Step 1}
\centering
\begin{tabular}{|l|l|l|}
\hline
Kill Chain Step & Attacking Action & Observation \\
\hline
Reconnaissance & nmap & Open ports, server versions \\
nikto & Directories and open ports &  \\
Weaponization and Delivery & Internet OSINT & Exploit on hand \\
\hline
\end{tabular}
\end{table}


\subsection{Cracking hashed passwords with `hashcat`}
Shortly after we got in, we headed right over to \texttt{/etc/passwd} to see what was there.  Surprised to see several applications with user accounts, like MySQL, somewhere in there we asked the computer a \texttt{ps -a} to show us all running processes since the databse port was not open on the other scan results.  Since the scans were limited to common ports with the settings we used, it could be that they were running MySQL with a custom port.  No dice.  Nothing was running.  With "x" in the field for passwords, we went to \texttt{/etc/shadow} and there was a file with encrypted passwords there.  



One password stood out:  the password for "Marlinspike," the administrator account on the target box.  Beginning with \texttt{\$6\$} and holding four \texttt{/} backslashes, I expected that this password might have been delineated.  A check with some resources suggested I was wrong. The \texttt{\$6\$} indicated a SHA-512 hash. \[[5]\]  Also, a closer examination of the shadow file showed us that many of those application accounts were basically null. \[[6]\]  There would be no way to log on to those.  

\begin{table}[H]
\caption{Encrypted Password}
\centering
\begin{tabular}{|l|l|l|l|l|l|l|l|l|l|l|l|l|l|l|l|l|l|l|l|l|l|l|l|l|l|l|l|l|l|l|l|l|}
\hline
Ciphertext & \$ & 6 & \$ & w & Q & b & 5 & n & V & 3 & T & \$ & x & B & Z & W & O & / &  &  &  &  &  &  &  &  &  &  &  &  &  &  \\
\hline
Notebook Notation &  & \# &  &  & \_ &  & \# &  & \_ & \# & \_ &  &  & \_ & \# & \_ & \_ &  &  &  &  &  &  &  &  &  &  &  &  &  &  &  \\
Ciphertext & j & O & k & b & n & 4 & t & 1 & R & U & I & L & r & c & k & w & 6 & 9 & L & R & / &  &  &  &  &  &  &  &  &  &  &  \\
Notebook Notation &  & \_ &  &  &  & \# &  & \# & \_ & \_ & \_ & \_ &  &  &  &  & \# & \# & \_ & \_ &  &  &  &  &  &  &  &  &  &  &  &  \\
Ciphertext & 0 & E & M & t & U & b & F & F & C & Y & p & M & 3 & M & U & H & V & m & t & y & Y & W & 9 & . & o & v & / &  &  &  &  &  \\
Notebook Notation & \# & \_ & \_ &  & \_ &  & \_ & \_ & \_ & \_ &  & \_ & \# & \_ & \_ & \_ & \_ &  &  &  & \_ & \_ & \# &  &  &  &  &  &  &  &  &  \\
Ciphertext & a & s & z & T & p & W & h & L & a & C & Z & x & 6 & F & v & y & 5 & t & p & U & U & x & Q & b & U & h & C & K & b & l & 4 & / \\
Notebook Notation &  &  &  & \_ &  & \_ &  & \_ &  & \_ & \# &  & \# & \_ &  &  & \# &  &  & \_ & \_ &  & \_ &  & \_ &  & \_ & \_ &  &  & \# &  \\
\hline
\end{tabular}
\end{table}

In the table above, we can see the password broken into four arbitrary sections, each ending in \texttt{/}.  Below each row of ciphertext is an example of a notation I use in my notebooks to help separate the numbers and the capital letters from the lowercase letters.  Since handwriting can be messy, using a row of notes below about each glyph can be helpful.  For example, a "5", an "S", and a lowercase "s" can all be easily confused when handwritten on a page. 

We wanted to crack that administrator password.  So, we set \texttt{hashcat} to the task.  Early on, \texttt{hashcat} would return "line length" exceptions.  This is a clue that we're asking \texttt{hashcat} to crack a hash that doesn't match the ciphertext that we've provided. \[[7]\]  A review of a Hashcat wiki page showed several examples that we could use to try to match our ciphertext with a type of hash. \[[9]\]  Since ours began with the distinctive \texttt{\$6\$}, had the context of being generated on a Unix system as part of \texttt{/etc/shadow}\[[5]\], and was the only kind on the Hashcat wiki with that feature:  our hash was SHA-512 (Unix).  This meant that our \texttt{-m} value needed to be 1800.  

\begin{lstlisting}
hashcat -m 1800 <TARGET HASH TO CRACK> /usr/share/wordlists/rockyou.txt --force
\end{lstlisting}

We boldly tried to crack this hash with Rock You in Kali.  We kept the computer running for several days.  Make sure the settings in Kali don't let the computer go to sleep after a few minutes.  As an extra precaution, consider using the Network settings on the VM hypervisor to disconnect the VM from the network altogether while it's running.  It's just not good to leave a computer running unattended for a long time if it can communicate with other machines.  


The result of our efforts:  nothing.  We didn't get the password.  Whatever it was, it wasn't in Rock You.  While there were other wordlists available in Kali, and we didn't even try to simple login guessing, it was time to move on to another task.  

\subsubsection{Analysis Check}
\begin{table}[H]
\caption{Discovery and Attack Cycle 2}
\centering
\begin{tabular}{|l|l|}
\hline
Action & Observation \\
\hline
Research & Identified hash name as SHA-512(Unix) \\
hashcat & Word list exhausted after 2 days and 16 hours. \\
\hline
\end{tabular}
\end{table}

\begin{table}[H]
\caption{Command to Kill Chain Step 2}
\centering
\begin{tabular}{|l|l|l|}
\hline
Kill Chain Step & Attacking Action & Observation \\
\hline
Weaponization & hashcat & Failed to crack hashed password. \\
\hline
\end{tabular}
\end{table}

\subsection{Roaming}
One of the first things we did, before we started "trying" to do anything, was roam around.  This informal reconnaissance proved to be quite instructive.  It let us know what was in the system, where we could go, and began to provide trickles of implementation-specific information about the target VM.

One of the first questions we asked when we popped these shells was, Where are we?  Metaspolit put us in the root directory upon launch, every time.


We plinked around the system, particularly near directories about passwords, logs, and the web servers.  We also explored some home directories of other users like marlinspike and root itself.  Included were checks on what kind of jobs the system might be running automatically.  



\subsubsection{Analysis Check}
\begin{table}[H]
\caption{Discovery and Attack Cycle 3}
\centering
\begin{tabular}{|l|l|}
\hline
Action & Observation \\
\hline
Target site navigation & Identified accounts, groups, critical programs \\
\texttt{ps -a}, crontabs, and application user accounts & showed what was and what was not running \\
\hline
\end{tabular}
\end{table}

\begin{table}[H]
\caption{Command to Kill Chain Step 3}
\centering
\begin{tabular}{|l|l|l|}
\hline
Kill Chain Step & Attacking Action & Observation \\
\hline
Reconnaissance & Common Unix commands & Identified useful targets in system \\
\hline
\end{tabular}
\end{table}

\subsection{The Ugly}
As part of describing our actions on the objective as "The Ugly, the Bad, and the Good," we'll open with the Ugly.  While roaming around the system as root with the shell created by the msfconsole exploit, we realized that there were already opportunities to begin covering tracks.  This wasn't as simple as it sounded to achieve the desired effect.  During this segment, we made a lot of mistakes.  Often, we had to recycle and re-attempt to get commands right.  If this had been a live pentest, there definitely would have been a lot of fumbling here; so, this block earns our title of "Ugly."

We roamed the system.  We received \texttt{/etc/passwd}, \texttt{crontab} files, \texttt{ps -a}, found a program called "popularity contest", and reviewed some Apache config scripts.  While hunting around in the Apache access log near \texttt{/var/log/apache2}, we saw direct evidence of our earlier \texttt{nikto} scans.  Each call that \texttt{nikto} made, while it wasn't shown to us on the attack box command line, that was recived by the Apache webserver was recorded in its logs.  In instances where those calls generated errors, usually because the requested page was not available, those calls were also recorded in the logs.  There was a lot of evidence just from our port scans that we had been interested in the target computer.


Right away, we set about deleting some of those logs.  Using simple commands, we were able to overwrite them.

\begin{lstlisting}
echo " " > access.log
\end{lstlisting}


However, when we would list the files, we could plainly see the time and date that we had modified them.  This just provided further evidence of our presence.  So, we began to investigate how those timestamps were used, where they came from, and what could be done about changing them.

Some simple attempts didn't work.  For example, merely changing the common time and date using ordinary UNIX commands would not override all of the fields.  Ideas like \texttt{date -s "2017-11-04 11:40:00" } and the like didn't work. \[[11]\] Our goal was to mimic the other timestamps that we saw on the box, which were made several months before, in November 2017. 


\begin{lstlisting}
stat access.log 
# revealed the three kinds of time related to a file.
touch -d "yyyy-mm-dd hh:mm:ss"
# could mimic the timestamps on a file.
\end{lstlisting}

Using some \texttt{touch} and \texttt{stat} commands, we were able to review the status of file times. It quickly became apparent that two out of three forms of file times could be readily changed with the change in system time, but one could not.  We needed to get time changed to a deeper level.  This period of experimentation, false starts, and failed tries is why we named this one "Ugly."  To keep the lessons clear, let's jump to what we learned:  in order to cover tracks effectively, it becomes very important to plan to control time on the targeted system.  \[[10]\] \[[11]\] \[[12]\] \[[13]\] \[[14]\] \[[15]\] \[[16]\] \[[17]\]


In retrospect, we think that it's likely that stopping time, changing time, and restarting time, might be among the more important activities in a phase of covering tracks.  This means that, for defensive actions, the routine and elementary monitoring of system time might be essential.  A log, off of the targeted system, recording what time a system believes it has, might be a valuable asset for pinpointing intrusions that involve an attacker altering file content in a stealthy way.  

One of the things we wanted to do as attackers was to not only destroy record of our activities, but to also make it seem like the actions had stopped with the last legitimate user's actions.  Stopping time was important for us here.  

A successful pattern of doing so might look like this, with \texttt{root}:
\begin{lstlisting}
# Stop coordination of time on a system
timedatectl set ntp off
# Observe the stop
timedatectl
# Set a time to a desired target time
timedatectl set-time "<TARGET_TIME>"
# Observe the time
timedatectl
# Conduct actions that will use the deceptive time
# Start time on the system clock
timedate set ntp on
# Observe the time
timedatectl
\end{lstlisting}

One of the other lessons learned was that the number of records that might need to be affected by these deceptive actions could be quite large.  In our targeted VM, we saw scripts that duplicated various logs automatically.  This means that not only might logs in the plain need to be altered, but also that zipped backup files stored elsewhere in the system might need it also.  Further, it was apparent that the more a sysadmin logged the system in various ways, the more logs there would need to be to contend with in this way.  \texttt{systemd} was still logging everything, including the time change commands.  Combined with common academic training in forensics that teaches us that nothing is ever really deleted:  well, covering tracks is unlikely to withstand any long-term, focused scrutiny.  At best, covering tracks might be camouflage and concealment, but not strong cover.


Perhaps we might grade and evaluate attackers based on how well they handle the "cover tracks" phase of hacking.  There are strong arguments for covering tracks:  doing it, not doing it, and maybe doing it well enough to get a little farther down the trail.  Let us just be aware of the pitfalls.  Also, it brings into focus the importance of OWASP's "Insufficient logging" as among the 2017 Top 10.  The more logging the better for the defender, in the attacker's phase of covering tracks.

\subsubsection{Analysis Check}
\begin{table}[H]
\caption{Discovery and Attack Cycle 4}
\centering
\begin{tabular}{|l|l|}
\hline
Action & Observation \\
\hline
Found Apache log files & Noticed evidence of previous intrusion \\
Destroyed files & Overwrote log files and observed timestamp change \\
Manipulated system clocks & Tinkered with clock times using various commands \\
Observed \texttt{systemd} logs & Witnessed some events of time changing on system \\
\hline
\end{tabular}
\end{table}

\begin{table}[H]
\caption{Command to Kill Chain Step 4}
\centering
\begin{tabular}{|l|l|l|}
\hline
Kill Chain Step & Attacking Action & Observation \\
\hline
Actions and Objectives & Log file destruction & Cosmetic covering of tracks \\
Installation & System clock manipulation & Supports installation of any file or follow on action by promoting camouflage of actions on objective \\
\hline
\end{tabular}
\end{table}

\section{The Bad}
We did some damaging on there.  As noted above, some log files were destroyed and their timestamps altered.  The attack didn't stop there.  We also:
- installed malicious user accounts with groups for \texttt{root sudo adm}
- tested and used \texttt{ssh} access for malicious user accounts
- added root back into the \texttt{adm} group for help with altering some files
- re-grouped the system's sole administrator out of the "adm" group
- deleted log backups 
- redirected users to another site 
- defaced the website.  


To accomplish all this, we used \texttt{vi} and the \texttt{os-shell} provided by the \texttt{msfconsole} exploit.  Most of the time, we used simple Unix commands to navigate to the right spot.  Then we would edit, delete, or create a file to meet out needs.  The PHP file to redirect users was a simple 5 line script, straight out of the PHP Man pages.  Editing user accounts and deleting files were common Unix admin tasks.  Damaging a system without regard in this manner was too easy.  We soon moved on to other tasks.

We would go on to attempt to install RTFM's "\texttt{ssh} Callback" script, and found some typos.  The goal of the script is to have the system call a known address at set time intervals.  This way, the compromised system would always phone home.

Later, in editing some \texttt{cron} jobs with \texttt{crontab} on old rule became apparent:  if you can't give the command in shell, then it won't work in script either.  So, it's generally better to give a command or two and then build up the script, troubleshooting along the way.  As with past cronjobs, I found myself triggering the jobs repeatedly to check to see if they worked as anticipated.  It's not \texttt{cron}'s fault; it's mine.  The Law of Shell Script Typos tells us that any syntax error that can muck up a cron job, will.  \[[27]\] \[[31]\] \[[32]\]

Just keep that in mind if you think you might type in a script on an unfamiliar system and hope it works on the first try.




\subsubsection{Analysis Check}
\begin{table}[H]
\caption{Discovery and Attack Cycle 5}
\centering
\begin{tabular}{|l|l|}
\hline
Action & Observation \\
\hline
Installed evil user account & Gained persistent access to target \\
Damaged legitimate administrator's group permissions & Expected damage to system authority \\
Defaced website & Denied access to expected information \\
Directed users to another site & DoS:  users completely denied web access to site through common redirection \\
\hline
\end{tabular}
\end{table}

\begin{table}[H]
\caption{Command to Kill Chain Step 5}
\centering
\begin{tabular}{|l|l|l|}
\hline
Kill Chain Step & Attacking Action & Observation \\
\hline
Installation & Installed malicious user account & Gained elevated permissions and persistent control \\
Command and Control & Damaged admin permissions for sole administrator & Changed balance of power over site to favor attacker \\
Actions and Objectives & Defaced website & Removed normal user ability to see the site they would ask for \\
Actions and Objectives & Redirected webserver traffic & Used simple PHP script to completely redirect traffic to google.com \\
Actions and Objectives & Installed ping script with cron job & Obscure web server file used to confirm contact with known ATKR IP \\
\hline
\end{tabular}
\end{table}

\section{The Good}
We began repairing the site, installing upgrades, and patching the vulns that got us in in the first place.  Our two big vulnerabilities were with ProFTPD and OpenSSH.  The first was a matter of applying a common upgrade.  The second required a little tinkering because Ubuntu was telling us we had the latest version, which wasn't quite true.

Using "eviluser" on \texttt{ssh}, we decided to patch ProFTPD with an upgrade, from 1.3.3x to 1.3.5a. \[[19]\]

\begin{lstlisting}
sudo dpkg --configure -a
sudo apt-get upgrade proftpd
\end{lstlisting}


System got many updates; 191 packages were upgraded; 126MB archives needed; had to add an additional 16.1MB of space on the disk.  Downloaded the update and did a reboot of the system.  While waiting for the installation process to complete, we noodled around on the Internet and found another exploit for ProFTPD. \[[21]\] \[[22]\] \[[26]\] One of the websites that mentioned it mentioned to also comment out a \texttt{mod\_copy.c} in its config files. \[[20]\] Downloaded 1.3.6 and proceeded to install that instead.

\begin{lstlisting}
sudo ./configure 
make install
\end{lstlisting}


Upgraded OpenSSH to 7.4p1 using a github gist from techguan.  Ubuntu's own measures seemed to suggest that the vuln 7.2 was the latest.  Once 7.4 was installed, we began checking on the exploits we had used to gain entry. We found that they were not working anymore. \[[23]\] \[[24]\] \[[33]\] \[[34]\]




\subsubsection{Analysis Check}
\begin{table}[H]
\caption{Discovery and Attack Cycle 6}
\centering
\begin{tabular}{|l|l|}
\hline
Action & Observation \\
\hline
Upgraded vulnerable programs & Removed vulnerabilities that got us in \\
\hline
\end{tabular}
\end{table}

\begin{table}[H]
\caption{Command to Kill Chain Step 6}
\centering
\begin{tabular}{|l|l|l|}
\hline
Kill Chain Step & Attacking Action & Observation \\
\hline
Installation & Installed ProFTPD upgrade & Closed exploit that got us in \\
Installation & Installed OpenSSH upgrade & Reduced attack potential \\
Exploitation & Logged on using evil user account after ssh upgrade & Confirmed account usage would persist after vuln closed \\
\hline
\end{tabular}
\end{table}

\section{Reporting}

\subsection{Follow-Up Vulnerability Scans}
After the repairs were done, we ran two automated scanners to see if they would detect any new problems.  Keep in mind, as these scans were being run, there was a known persistent malicious account on the box.  

The first scan we ran was an OWASP ZAP, which primarily tests web applications.  We used the "quick scan" feature.  It found no significant vulnerablilities, save for an alleged RFI vuln on "google.com".  We were redirecting users there.  Let us recognize that by destroying most of the website, there simply wasn't any mechanism to get in.  


Next, we conducted a simple Nessus Basic Network scan.  While Nessus did tell us that the ports were open for FTP, SSH, and HTTP, it did not tell us about the performance of those related programs.  For example, we didn't see evidence that our programs serving those applications were vulnerable.  However, notice we had a malicious user account installed on the \texttt{ssh} with a weak password.  This was not detected.


So, one of the things we see in action here is that a vulnerability may not seem to be present, but that doesn't exclude any well-used weaknesses.  That is, some activities may simply be out of bounds of the scan.

\subsection{Lessons learned from Basic Pentesting 1}

We learned a lot from this exercise.  In the first few discovery cycles, we could see that familiarity with some basic scans, combined with some quick research, allowed us to dredge up some leads to quickly exploit the target with known metasploit modules.  Once inside, we learned about what past admins had done by observing what accounts had been created and what processes were active.  We could see that scans, like most other calls to web servers, were going to leave their own traces in the logs.  When it came time to change the time to try to cover tracks, we could see that time changes themselves would leave evidence.  Also, we can see that despite this, changing the system time and hardware clocks allowed us to be capable of concealing multiple activities during the time that the clock was changed in the attacker's favor.

As we began to directly damage and repair the system, we could see it was very easy to destroy files and change configuration.  As expected, root and some simple Unix commands were all that was needed to completely take over the machine.  With the addition of our own user account, we were able to persist access after the updates were done to close the vulnerabilities that got us in.  As we wrapped up our repairs, we remembered to perform a check on our actions.  By opening and closing the exercise with vulnerability scans, we were able to produce quantifiable and qualifiable results that showed the change in the system's status as a result of our activities.  

If our purpose for pentesting is to provide a mechanism for repair, then in this instance we proved some repairs could be made.  As we could see by not being able to re-exploit those old vulns, and doing the research to skip over some others, we were able to improve the overall fortification of the system.  All in, it was a good exercise for us.

\section{Annotated Bibliography}


\endnote{[1]}  Brotherston, Lee and Amanda Berlin.  Defensive Security Handbook:  Best Practices for Securing Infrastructure.  O'Reilly, April 2017. pp. 6-9. ISBN 978-1-491-96038-7. 

A description of Lockheed-Martin's Intusion Kill Chain.  Brotherston and Berlin applied their use case of an overview of a ransomware attack to the phases of Lockheed's kill chain.  Their definitions and example serves as guidance for the contstruction of our own, here.


\endnote{[2]}  INTERNET: \url{https://nvlpubs.nist.gov/nistpubs/Legacy/SP/nistspecialpublication800-115.pdf}\endnote{\url{https://nvlpubs.nist.gov/nistpubs/Legacy/SP/nistspecialpublication800-115.pdf}}
NIST documents includes descriptions of vulnerability scanning, penetration testing, password cracking, and other critical terms of applicable techniques.

\endnote{[3]}  Pierce, Josiah.  \url{https://www.vulnhub.com/entry/basic-pentesting-1,216/}\endnote{\url{https://www.vulnhub.com/entry/basic-pentesting-1,216/}}
VulnHub virtual machine target for pentesting practice.

\[[4]\] INTERNET:  \url{https://www.rapid7.com/db/modules/exploit/unix/ftp/proftpd\_133c\_backdoor}\endnote{\url{https://www.rapid7.com/db/modules/exploit/unix/ftp/proftpd_133c_backdoor}}
Rapid7's page on an exploit for ProFTPD 1.3.3c.

\[[5]\] INTERNET:  \url{https://unix.stackexchange.com/questions/8229/what-methods-are-used-to-encrypt-passwords-in-etc-passwd-and-etc-shadow}\endnote{\url{https://unix.stackexchange.com/questions/8229/what-methods-are-used-to-encrypt-passwords-in-etc-passwd-and-etc-shadow}}
Question told us about \texttt{\$6\$} use in \texttt{/etc/shadow} as an indication of an encryption method.  This notation says it's SHA-512.

\[[6]\] INTERNET:  \url{https://unix.stackexchange.com/questions/252016/difference-between-vs-vs-in-etc-shadow}\endnote{\url{https://unix.stackexchange.com/questions/252016/difference-between-vs-vs-in-etc-shadow}}
Question gave us clues to understanding the use of exclamation points and asterisks in \texttt{etc/shadow}.

\[[7]\] INTERNET:  \url{https://www.unix-ninja.com/p/Hashcat\_Line\_Length\_Exceptions}\endnote{\url{https://www.unix-ninja.com/p/Hashcat_Line_Length_Exceptions}}
Blog post pointing out that line length exceptions are often related to misconfiguring the command for the type of ciphertext provided.

\[[8]\] INTERNET:  \url{https://hashcat.net/forum/thread-3399.html}\endnote{\url{https://hashcat.net/forum/thread-3399.html}}
Hashcat forum thread that helped us find the Hashcat wiki with the examples.

\[[9]\] INTERNET:  \url{https://hashcat.net/wiki/doku.php?id=example\_hashes}\endnote{\url{https://hashcat.net/wiki/doku.php?id=example_hashes}}
A listing of example hashes and the ciphers that generate them.

\[[10]\] INTERNET:  \url{http://www.informit.com/articles/article.aspx?p=1161217\&seqNum=11}\endnote{\url{http://www.informit.com/articles/article.aspx?p=1161217&seqNum=11}}
Using \texttt{hwclock} commands for adjusting the hardware clock.

\[[11]\]  INTERNET:  \url{https://www.garron.me/en/linux/set-time-date-timezone-ntp-linux-shell-gnome-command-line.html}\endnote{\url{https://www.garron.me/en/linux/set-time-date-timezone-ntp-linux-shell-gnome-command-line.html}}
Setting date and time on a UNIX system.

\[[12]\]  INTERNET:  \url{https://www.thegeekstuff.com/2012/11/linux-touch-command/?utm\_source=feedburner}\endnote{\url{https://www.thegeekstuff.com/2012/11/linux-touch-command/?utm_source=feedburner}}
Using \texttt{touch} and \texttt{stat}, with examples.

\[[13]\]  Hudson, Andrew and Paul Hudson.  INTERNET: \url{http://www.informit.com/articles/article.aspx?p=1161217\&seqNum=11}\endnote{\url{http://www.informit.com/articles/article.aspx?p=1161217&seqNum=11}}
Online book chapter explaining date and time setting changes as they are used in Ubuntu Linux.

\[[14]\]  Garron, Guillermo.  "Set Time, Date Timezone in Linux from Command Line or Gnome | Use ntp" INTERNET: \url{https://www.garron.me/en/linux/}\endnote{\url{https://www.garron.me/en/linux/}} 2012.
Setting date and time in Ubuntu, including \texttt{hwclock} commands.

\[[15]\]  INTERNET: \url{https://stackoverflow.com/questions/16126992/setting-changing-the-ctime-or-change-time-attribute-on-a-file/17066309#17066309}% \endnote{\protect\url{https://stackoverflow.com/questions/16126992/setting-changing-the-ctime-or-change-time-attribute-on-a-file/17066309#17066309}}
Forum answer which turned us towards the idea that there were multiple kinds of time in Linux systems, \texttt{ctime} and \texttt{mtime}.

\[[16]\]  Ganapathy, Lakshmanan. "5 Linux Touch Command Examples (How to Change File Timestamp)." INTERNET: \url{https://www.thegeekstuff.com/2012/11/linux-touch-command/?utm\_source=feedburner}\endnote{\url{https://www.thegeekstuff.com/2012/11/linux-touch-command/?utm_source=feedburner}} 2012.
Examples of using \texttt{touch} and \texttt{stat} to read and write file attributes.

\[[17]\]  INTERNET: \url{https://askubuntu.com/questions/920598/disable-network-time-synchronization-in-ubuntu-16-04}\endnote{\url{https://askubuntu.com/questions/920598/disable-network-time-synchronization-in-ubuntu-16-04}}
Forum post that clued us in on using commands like \texttt{timedatectl}, with examples.

\[[18]\]  INTERNET: \url{https://askubuntu.com/questions/24006/how-do-i-reset-a-lost-administrative-password}\endnote{\url{https://askubuntu.com/questions/24006/how-do-i-reset-a-lost-administrative-password}}
Forum post with hyperlinks to documentation; this gave us leads on how much trouble it might be to change the root or administrator's password without knowing the original.

\[[19]\]  INTERNET: \url{http://www.proftpd.org/docs/howto/Upgrade.html}\endnote{\url{http://www.proftpd.org/docs/howto/Upgrade.html}}
ProFTPD official directions on how to upgrade their program.

\[[20]\]  INTERNET: \url{https://www.virtualmin.com/node/43086}\endnote{\url{https://www.virtualmin.com/node/43086}}
Forum post answer on how to manually \texttt{wget} a copy of an upgrade to ProFTPD and install it.  Clues that there were also problems and vulnerabilities with \texttt{mod\_copy.c} modules of the program.

\[[21]\]  INTERNET: \url{http://www.proftpd.org/docs/contrib/mod\_copy.html}\endnote{\url{http://www.proftpd.org/docs/contrib/mod_copy.html}}
ProFTPD documentation on what \texttt{mod\_copy.c} does in their programs.

\[[22]\]  INTERNET: \url{https://www.rapid7.com/db/modules/exploit/unix/ftp/proftpd\_modcopy\_exec}\endnote{\url{https://www.rapid7.com/db/modules/exploit/unix/ftp/proftpd_modcopy_exec}}
Rapid7 page on the metasploit module to exploit a vuln in \texttt{mod\_copy.c}.  Exploit name is \texttt{proftpd\_modcopy\_exec}.

\[[23]\]  INTERNET: \url{https://gist.github.com/stefansundin/0fd6e9de172041817d0b8a75f1ede677}\endnote{\url{https://gist.github.com/stefansundin/0fd6e9de172041817d0b8a75f1ede677}}
Github Gist that suggested it was possible to get a 7.3 OpenSSH in Ubuntu.  At the time we did our upgrades the normal way, only a vulnerable 7.2 version seemed to be available with those techniques.

\[[24]\]  INTERNET: \url{https://gist.github.com/techgaun/df66d37379df37838482c4c3470bc48e}\endnote{\url{https://gist.github.com/techgaun/df66d37379df37838482c4c3470bc48e}}
Github Gist that explained how to install OpenSSH version 7.4 in Ubuntu.  These directions allowed us to clear some known vulnerabilities in OpenSSH.

\[[25]\]  "Ubuntu Linux: Start / Stop / Restart / Reload OpenSSH Server"  INTERNET: \url{https://www.cyberciti.biz/faq/howto-start-stop-ssh-server/}\endnote{\url{https://www.cyberciti.biz/faq/howto-start-stop-ssh-server/}}
Blog post with examples about starting and stopping an ssh service.

\[[26]\]  INTERNET: \url{https://esc.sh/blog/proftp-vulnerability-could-allow-an-attacker-to-gain-a-shell-in-your-server/}\endnote{\url{https://esc.sh/blog/proftp-vulnerability-could-allow-an-attacker-to-gain-a-shell-in-your-server/}}
Blog post showing how to use metasploit to exploit the \texttt{mod\_copy} vulnerability in ProFTPD.

\[[27]\]  "Linux: List / Display All Cron Jobs"  INTERNET: \url{https://www.cyberciti.biz/faq/linux-show-what-cron-jobs-are-setup/}\endnote{\url{https://www.cyberciti.biz/faq/linux-show-what-cron-jobs-are-setup/}}
A review of commands used for storing and editing \texttt{cron} jobs in Linux.  This model uses a breakdown of four classes of jobs:  hourly, daily, weekly, and monthly jobs are all stored separately. 

\[[28]\]  INTERNET: \url{https://null-byte.wonderhowto.com/how-to/create-reverse-shell-remotely-execute-root-commands-over-any-open-port-using-netcat-bash-0132658/}\endnote{\url{https://null-byte.wonderhowto.com/how-to/create-reverse-shell-remotely-execute-root-commands-over-any-open-port-using-netcat-bash-0132658/}}
A refresher on the basics of a \texttt{nc} netcat reverse shell.

\[[29]\]  Li, Zhaoyu.  "A Simple Bash Script to Reverse SSH Tunnel Automatically"  INTERNET: \url{https://zhaoyu.li/post/auto-reverse-ssh/}\endnote{\url{https://zhaoyu.li/post/auto-reverse-ssh/}}
Example of a brief reverse shell script that tries to reconnect automatically every ten seconds.

\[[30]\]  Gargiullo, Michael.  "Fun With SSH Reverse Shells"  INTERNET: \url{https://www.pivotpointsecurity.com/blog/fun-with-ssh-reverse-shells/}\endnote{\url{https://www.pivotpointsecurity.com/blog/fun-with-ssh-reverse-shells/}}
Blog post that reminded us to use commands like \texttt{ssh -p <PORTNUMBER>} to ssh back into a specific port.

\[[31]\]  "Linux Start Restart and Stop The Cron or Crond Service"  INTERNET: \url{https://www.cyberciti.biz/faq/howto-linux-unix-start-restart-cron/}\endnote{\url{https://www.cyberciti.biz/faq/howto-linux-unix-start-restart-cron/}} 2017.
Example commands on restarting \texttt{cron}.

\[[32]\]  INTERNET: \url{https://askubuntu.com/questions/2368/how-do-i-set-up-a-cron-job}\endnote{\url{https://askubuntu.com/questions/2368/how-do-i-set-up-a-cron-job}}
Blog post showing the setup of cron jobs in environments where crontab might file the jobs under the separate folders for: hourly, dailly, weekly and monthly.

\[[33]\]  INTERNET: \url{https://ubuntuforums.org/showthread.php?t=2380028}\endnote{\url{https://ubuntuforums.org/showthread.php?t=2380028}}
Forum post on upgrading OpenSSH on Ubuntu.

\[[34]\]  INTERNET: \url{https://stackoverflow.com/questions/36453976/upgrade-openssh-7-2p-in-ubuntu-14-04}\endnote{\url{https://stackoverflow.com/questions/36453976/upgrade-openssh-7-2p-in-ubuntu-14-04}}
Forum post that showed an example of upgrading to a specific version of OpenSSH.  Slight modifications on these commands allowed us to get the version we wanted installed on the target box.

\endnote{[35]}  Clark, Ben. RTFM:  Red Team Field Manual.  ISBN 9781494295509.
A cookbook of Red Team scripts for field use.  Callback SSH script referenced in this post was listed on p. 9.

References
[1]:  ``
[2]: \url{https://nvlpubs.nist.gov/nistpubs/Legacy/SP/nistspecialpublication800-115.pdf}
[3]: \url{https://www.vulnhub.com/entry/basic-pentesting-1,216/}
[4]: \url{https://www.rapid7.com/db/modules/exploit/unix/ftp/proftpd_133c_backdoor}
[5]: \url{https://unix.stackexchange.com/questions/8229/what-methods-are-used-to-encrypt-passwords-in-etc-passwd-and-etc-shadow}
[6]: \url{https://unix.stackexchange.com/questions/252016/difference-between-vs-vs-in-etc-shadow}
[7]: \url{https://www.unix-ninja.com/p/Hashcat_Line_Length_Exceptions}
[8]: \url{https://hashcat.net/forum/thread-3399.html}
[9]: \url{https://hashcat.net/wiki/doku.php?id=example_hashes}
[10]: \url{http://www.informit.com/articles/article.aspx?p=1161217&seqNum=11}
[11]: \url{https://www.garron.me/en/linux/set-time-date-timezone-ntp-linux-shell-gnome-command-line.html}
[12]: \url{https://www.thegeekstuff.com/2012/11/linux-touch-command/?utm_source=feedburner}
[13]: \url{http://www.informit.com/articles/article.aspx?p=1161217&seqNum=11}
[14]: \url{https://www.garron.me/en/linux/set-time-date-timezone-ntp-linux-shell-gnome-command-line.html}
[15]: \url{https://stackoverflow.com/questions/16126992/setting-changing-the-ctime-or-change-time-attribute-on-a-file/17066309#17066309}
[16]: \url{https://www.thegeekstuff.com/2012/11/linux-touch-command/?utm_source=feedburner}
[17]: \url{https://askubuntu.com/questions/920598/disable-network-time-synchronization-in-ubuntu-16-04}
[18]: \url{https://askubuntu.com/questions/24006/how-do-i-reset-a-lost-administrative-password}
[19]: \url{http://www.proftpd.org/docs/howto/Upgrade.html}
[20]: \url{https://www.virtualmin.com/node/43086}
[21]: \url{http://www.proftpd.org/docs/contrib/mod_copy.html}
[22]: \url{https://www.rapid7.com/db/modules/exploit/unix/ftp/proftpd_modcopy_exec}
[23]: \url{https://gist.github.com/stefansundin/0fd6e9de172041817d0b8a75f1ede677}
[24]: \url{https://gist.github.com/techgaun/df66d37379df37838482c4c3470bc48e}
[25]: \url{https://www.cyberciti.biz/faq/howto-start-stop-ssh-server/}
[26]: \url{https://esc.sh/blog/proftp-vulnerability-could-allow-an-attacker-to-gain-a-shell-in-your-server/}
[27]: \url{https://www.cyberciti.biz/faq/linux-show-what-cron-jobs-are-setup/}
[28]: \url{https://null-byte.wonderhowto.com/how-to/create-reverse-shell-remotely-execute-root-commands-over-any-open-port-using-netcat-bash-0132658/}
[29]: \url{https://zhaoyu.li/post/auto-reverse-ssh/}
[30]: \url{https://www.pivotpointsecurity.com/blog/fun-with-ssh-reverse-shells/}
[31]: \url{https://www.cyberciti.biz/faq/howto-linux-unix-start-restart-cron/}
[32]: \url{https://askubuntu.com/questions/2368/how-do-i-set-up-a-cron-job}
[33]: \url{https://ubuntuforums.org/showthread.php?t=2380028}
[34]: \url{https://stackoverflow.com/questions/36453976/upgrade-openssh-7-2p-in-ubuntu-14-04}


\markchapterendnotes
